\chapter{Philosophy}
\section{Scope}
Here I will write general notes from sources on the subject of philosophy.
\section{Nicomachean Ethics}
\subsection{3/26/26}
While I have constantly defended virtue ethics, and even Aristotle's ethics, I am a bit critical of his methods. Rather than defining a platonic ideal, he observes those around him and defines ethics based upon that. It works quite well and is something I also do but this as the primarly method tends to lead to some differences.

I am also surprised how similar his views of ethics are so similar to 21st century ethics in many ways; considering how different the views of people just a couple decades ago were.

I definitely find coherence with his idea of magnanimity with my philotimic virtue. While I had originally conceived of the idea from listening to a summary of the idea, reading it directly from him shows that it was incredibly similar. Making me almost wonder if I should forget having my own word and just use that one. The only real difference is that I also blend philotimic with my specific views of morality and the idea of self-authorship; obviously an Aristotelian magnamonious man would value self-authorship, it isn't directly imbued with the word. Also, his spesific definition has a more prideful connotation.

Kinda funny how even in that time older folks found young comedians as vulger and inappropriate while the younger found the otherside as offensive and cruel.

Obviously his views of justice, while different are quite well.
